% GENERATED FILE. Edit canonical lessons, then run npm run generate:books.
% canonical-source-hash: fnv1a64:1ceb02e46ecd333f
% canonical-lessons: SA-C315-visphotah, SA-C315-pratidhvanih, SA-C315-akasmikah, SA-C315-udvijate, SA-C315-tvarate

\chapter{A Blast}
\label{ch:sa-a-blast}
\hlchaptermodality{315}

\begin{chapteropening}
\textbf{By the end of this chapter:} \emph{I can describe a sudden loud event.}

Complete the last lesson of chapter 315: I can describe a sudden loud event.
\end{chapteropening}

\section[\texorpdfstring{\sk{विस्फोटः} (visphoṭaḥ)}{विस्फोटः (visphoṭaḥ)}]{\sk{विस्फोटः} (visphoṭaḥ) — an explosion}

\label{lesson:SA-C315-visphotah}



\begin{quote}
Before the new one: say the Sanskrit for honour, then the Sanskrit for congratulates.
\end{quote}

\subsection*{You'll want to know: \sk{विस्फोटः}}
\textbf{\sk{विस्फोटः}} — \emph{visphoṭaḥ} — \textquotedblleft{}an explosion\textquotedblright{}.

\begin{quote}
\sk{प्रातः} \sk{नगरे} \sk{विस्फोटः} \sk{अभवत्।}
\end{quote}

\begin{quote}
\emph{prātaḥ nagare visphoṭaḥ abhavat.}
\end{quote}

\begin{quote}
There was an explosion in the city in the morning.
\end{quote}

A masculine noun.

\begin{cousinweb}
\textbf{\sk{विस्फोटः}} (\emph{visphoṭaḥ}) is built from \emph{vi} and the root √sphuṭ, \textquotedblleft{}to burst\textquotedblright{}, as in \emph{sphuṭati}; MW gives \emph{visphoṭa} as cracking, crashing, and an explosion is its modern use (MW).
\end{cousinweb}

\begin{grammarlens}[title={What you've built}]
One more word for telling what happened.
\end{grammarlens}

\subsection*{Your turn}
\begin{itemize}
\raggedright
  \item \emph{Say it:} \emph{visphoṭaḥ}
  \item \emph{Say it:} \emph{visphoṭaḥ}, once more
  \item \emph{Say it:} \emph{prātaḥ nagare visphoṭaḥ abhavat.} — There was an explosion in the city in the morning.
  \item \emph{Recall:} say the Sanskrit for honour, then the Sanskrit for congratulates, then say \emph{visphoṭaḥ} again
\end{itemize}

\begin{tcolorbox}[breakable,colback=teal!4,colframe=teal!35!black,title={Before you move on}]
What does \sk{विस्फोटः} mean? (\textquotedblleft{}An explosion\textquotedblright{}.) Say it once more.
\end{tcolorbox}
\section[\texorpdfstring{\sk{प्रतिध्वनिः} (pratidhvaniḥ)}{प्रतिध्वनिः (pratidhvaniḥ)}]{\sk{प्रतिध्वनिः} (pratidhvaniḥ) — an echo}

\label{lesson:SA-C315-pratidhvanih}



\begin{quote}
Before the new one: say the Sanskrit for congratulates, then the Sanskrit for an explosion.
\end{quote}

\subsection*{You'll want to know: \sk{प्रतिध्वनिः}}
\textbf{\sk{प्रतिध्वनिः}} — \emph{pratidhvaniḥ} — \textquotedblleft{}an echo\textquotedblright{}.

\begin{quote}
\sk{पर्वतेषु} \sk{विस्फोटस्य} \sk{प्रतिध्वनिः} \sk{भवति।}
\end{quote}

\begin{quote}
\emph{parvateṣu visphoṭasya pratidhvaniḥ bhavati.}
\end{quote}

\begin{quote}
The blast echoes in the hills.
\end{quote}

A masculine noun.

\begin{cousinweb}
\textbf{\sk{प्रतिध्वनिः}} (\emph{pratidhvaniḥ}) is built from \emph{prati}, \textquotedblleft{}back, against\textquotedblright{}, and \emph{dhvani}, \textquotedblleft{}a sound\textquotedblright{}, from the root √dhvan; MW gives \emph{pratidhvani} as an echo (MW).
\end{cousinweb}

\begin{grammarlens}[title={What you've built}]
One more word for telling what happened.
\end{grammarlens}

\subsection*{Your turn}
\begin{itemize}
\raggedright
  \item \emph{Say it:} \emph{pratidhvaniḥ}
  \item \emph{Say it:} \emph{pratidhvaniḥ}, once more
  \item \emph{Say it:} \emph{parvateṣu visphoṭasya pratidhvaniḥ bhavati.} — The blast echoes in the hills.
  \item \emph{Recall:} say the Sanskrit for congratulates, then the Sanskrit for an explosion, then say \emph{pratidhvaniḥ} again
\end{itemize}

\begin{tcolorbox}[breakable,colback=teal!4,colframe=teal!35!black,title={Before you move on}]
What does \sk{प्रतिध्वनिः} mean? (\textquotedblleft{}An echo\textquotedblright{}.) Say it once more.
\end{tcolorbox}
\section[\texorpdfstring{\sk{आकस्मिकः} (ākasmikaḥ)}{आकस्मिकः (ākasmikaḥ)}]{\sk{आकस्मिकः} (ākasmikaḥ) — sudden; unexpected}

\label{lesson:SA-C315-akasmikah}



\begin{quote}
Before the new one: say the Sanskrit for an explosion, then the Sanskrit for an echo.
\end{quote}

\subsection*{You'll want to know: \sk{आकस्मिकः}}
\textbf{\sk{आकस्मिकः}} — \emph{ākasmikaḥ} — \textquotedblleft{}sudden; unexpected\textquotedblright{}.

\begin{quote}
\sk{आकस्मिकात्} \sk{शब्दात्} \sk{बालाः} \sk{भीताः।}
\end{quote}

\begin{quote}
\emph{ākasmikāt śabdāt bālāḥ bhītāḥ.}
\end{quote}

\begin{quote}
The children are frightened by the sudden noise.
\end{quote}

\begin{cousinweb}
\textbf{\sk{आकस्मिकः}} (\emph{ākasmikaḥ}) is built on \emph{akasmāt}, \textquotedblleft{}without a cause, suddenly\textquotedblright{} (from \emph{a}, \textquotedblleft{}not\textquotedblright{}, and \emph{kasmāt}, \textquotedblleft{}from what?\textquotedblright{}); MW gives \emph{ākasmika} as causeless, unforeseen, sudden (MW).
\end{cousinweb}

\begin{grammarlens}[title={What you've built}]
One more word for telling what happened.
\end{grammarlens}

\subsection*{Your turn}
\begin{itemize}
\raggedright
  \item \emph{Say it:} \emph{ākasmikaḥ}
  \item \emph{Say it:} \emph{ākasmikaḥ}, once more
  \item \emph{Say it:} \emph{ākasmikāt śabdāt bālāḥ bhītāḥ.} — The children are frightened by the sudden noise.
  \item \emph{Recall:} say the Sanskrit for an explosion, then the Sanskrit for an echo, then say \emph{ākasmikaḥ} again
\end{itemize}

\begin{tcolorbox}[breakable,colback=teal!4,colframe=teal!35!black,title={Before you move on}]
What does \sk{आकस्मिकः} mean? (\textquotedblleft{}Sudden; unexpected\textquotedblright{}.) Say it once more.
\end{tcolorbox}
\section[\texorpdfstring{\sk{उद्विजते} (udvijate)}{उद्विजते (udvijate)}]{\sk{उद्विजते} (udvijate) — is alarmed; starts in fear}

\label{lesson:SA-C315-udvijate}



\begin{quote}
Before the new one: say the Sanskrit for an echo, then the Sanskrit for unexpected.
\end{quote}

\subsection*{You'll want to know: \sk{उद्विजते}}
\textbf{\sk{उद्विजते}} — \emph{udvijate} — \textquotedblleft{}is alarmed; starts in fear\textquotedblright{}.

\begin{quote}
\sk{शब्दात्} \sk{सर्वे} \sk{उद्विजन्ते।}
\end{quote}

\begin{quote}
\emph{śabdāt sarve udvijante.}
\end{quote}

\begin{quote}
Everyone is alarmed by the noise.
\end{quote}

Past (he …): \sk{सः} \sk{शब्दात्} \sk{उद्विग्नवान्}, \emph{saḥ śabdāt udvignavān}.

\begin{cousinweb}
\textbf{\sk{उद्विजते}} (\emph{udvijate}) is built from \emph{ud} and the root √vij, \textquotedblleft{}to start, move with a jerk\textquotedblright{}, as in \emph{vega}; MW gives \emph{ud-vij} as to shudder, tremble, start and fear (MW).
\end{cousinweb}

\begin{grammarlens}[title={What you've built}]
One more word for telling what happened.
\end{grammarlens}

\subsection*{Your turn}
\begin{itemize}
\raggedright
  \item \emph{Say it:} \emph{udvijate}
  \item \emph{Say it:} \emph{udvijate}, once more
  \item \emph{Say it:} \emph{śabdāt sarve udvijante.} — Everyone is alarmed by the noise.
  \item \emph{Recall:} say the Sanskrit for an echo, then the Sanskrit for unexpected, then say \emph{udvijate} again
\end{itemize}

\begin{tcolorbox}[breakable,colback=teal!4,colframe=teal!35!black,title={Before you move on}]
What does \sk{उद्विजते} mean? (\textquotedblleft{}Is alarmed; starts in fear\textquotedblright{}.) Say it once more.
\end{tcolorbox}
\section[\texorpdfstring{\sk{त्वरते} (tvarate)}{त्वरते (tvarate)}]{\sk{त्वरते} (tvarate) — hurries}

\label{lesson:SA-C315-tvarate}



\begin{quote}
Before the new one: say the Sanskrit for unexpected, then the Sanskrit for is alarmed.
\end{quote}

\subsection*{You'll want to know: \sk{त्वरते}}
\textbf{\sk{त्वरते}} — \emph{tvarate} — \textquotedblleft{}hurries\textquotedblright{}.

\begin{quote}
\sk{सर्वे} \sk{चिकित्सालयं} \sk{प्रति} \sk{त्वरन्ते।}
\end{quote}

\begin{quote}
\emph{sarve cikitsālayaṃ prati tvarante.}
\end{quote}

\begin{quote}
Everyone hurries towards the hospital.
\end{quote}

Past (he …): \sk{सः} \sk{चिकित्सालयं} \sk{प्रति} \sk{त्वरितवान्}, \emph{saḥ cikitsālayaṃ prati tvaritavān}.

\begin{cousinweb}
\textbf{\sk{त्वरते}} (\emph{tvarate}) comes from the root √tvar, \textquotedblleft{}to hurry, make haste\textquotedblright{} (MW).
\end{cousinweb}

\begin{grammarlens}[title={What you've built}]
One more word for telling what happened.
\end{grammarlens}

\subsection*{Your turn}
\begin{itemize}
\raggedright
  \item \emph{Say it:} \emph{tvarate}
  \item \emph{Say it:} \emph{tvarate}, once more
  \item \emph{Say it:} \emph{sarve cikitsālayaṃ prati tvarante.} — Everyone hurries towards the hospital.
  \item \emph{Recall:} say the Sanskrit for unexpected, then the Sanskrit for is alarmed, then say \emph{tvarate} again
\end{itemize}

\begin{tcolorbox}[breakable,colback=teal!4,colframe=teal!35!black,title={Before you move on}]
What does \sk{त्वरते} mean? (\textquotedblleft{}Hurries\textquotedblright{}.) Say it once more.
\end{tcolorbox}
